% =============================================================================
% Бүлэг 6: Үр дүн ба хэлэлцүүлэг
% =============================================================================

\chapter{Үр дүн ба хэлэлцүүлэг}
\label{ch:results}
\pagecolor{white}

Бүх дүн кирилл асуулт дээр. Тэмдэглэгээ: \textbf{T4} --- Kaggle T4, fp16, LoRA-г нэгтгэсэн жин; \textbf{A100} --- Vast
A100, bf16, LoRA adapter; \textbf{GPU-гүй} --- өмнөх run-ийн хувилбаруудаас загваргүй дахин сонголт.

\section{SQL үүсгэх загварын үр дүн}
\label{sec:numbers}

\textbf{Dev (300 = 224 / 38 / 38).} Хүснэгт~\ref{tab:dev-full}-д задаргаатай үр дүнг, Зураг~\ref{fig:ladder}-т ахицын
замыг үзүүлэв.

\begin{table}[H]
\centering\small
\caption{Dev-ийн үр дүн (задаргаатай)}
\label{tab:dev-full}
\begin{adjustbox}{max width=\textwidth}
\begin{tabular}{llrrrrl}
\toprule
ID & Тохиргоо & Нийт /300 & Хариулах /224 & Боломжгүй /38 & Хоёрдмол /38 & GPU \\
\midrule
D2 & LoRA + RAG & 47 (15.7\%) & 37 (16.5\%) & 0 & 10 & A100 \\
D2 tr & + англи орчуулга & 81 (27.0\%) & 73 (32.6\%) & 0 & 8 & A100 \\
H1 & D2 + засвар, татгалзал & 74 (24.7\%) & 47 (21.0\%) & 17 & 10 & A100 \\
H4 & H1 + загварт англи, хайлт монголоор & 110 (36.7\%) & 82 (36.6\%) & 16 & 12 & A100 \\
E1 & H4 + алдаагаар засах & 111 (37.0\%) & 86 (38.4\%) & 14 & 11 & A100 \\
E2 & E1 + 8 хувилбар, санал & \textbf{123 (41.0\%)} & \textbf{100 (44.6\%)} & 13 & 10 & A100 \\
E3c & E2 + тодруулга (итгэл $<$ 0.5) & 137 (45.7\%) & 88 (39.3\%) & 30 & 19 & GPU-гүй \\
E4 & E2 + шүүлт + dev-ийн дүрэм & 159 (53.0\%) & 107 (47.8\%) & --- & --- & GPU-гүй \\
I3t & H4 + 579 жишээ + засварласан орчуулга & 132 (44.0\%) & 105 (46.9\%) & --- & --- & T4 \\
M1a & I3 + сургасан MADLAD орчуулга (M1) & 160 (53.3\%) & \textbf{127 (56.7\%)} & --- & --- & T4 \\
\bottomrule
\end{tabular}
\end{adjustbox}
\par\smallskip\footnotesize Хариулах боломжтой EX-ийн 95\% итгэх завсар: D2 11.6--21.4\%, H4 30.4--42.9\%, E2 38.4--51.3\%,
M1a 50.5--63.4\%.
\end{table}

\begin{figure}[htbp]
\centering
\begin{tikzpicture}
\begin{axis}[
  ybar, bar width=14pt, width=\textwidth, height=5.6cm,
  ymin=0, ymax=60, ylabel={Нийт зөв, \%}, ylabel style={font=\small},
  symbolic x coords={A7,B5,D2,D2tr,H4,E2,E3c,E4},
  xticklabels={A7,B5,D2,D2 tr,H4,E2,E3c,E4},
  xtick=data, ticklabel style={font=\small},
  nodes near coords, nodes near coords style={font=\scriptsize, /pgf/number format/.cd, fixed, fixed zerofill, precision=1},
  ymajorgrids, grid style={gray!20}, axis x line*=bottom, axis y line*=left]
\addplot[fill=orange!45, draw=orange!70!black] coordinates {(A7,2.3) (B5,14.3) (D2,15.7) (D2tr,27.0) (H4,36.7) (E2,41.0) (E3c,45.7) (E4,53.0)};
\end{axis}
\end{tikzpicture}
\caption{Dev дээрх ахицын зам (нийт зөв, \%)}
\label{fig:ladder}
\end{figure}

\textbf{Test (200 = 151 / 24 / 25).}\label{sec:test-results} Test-ийн асуултууд dev-ээс өөр загварын гэр бүлд хамаардаг
тул тохиргоо шинэ асуултад хэр ерөнхийлж байгааг харуулна (Хүснэгт~\ref{tab:test}); статистик ач холбогдлыг
Хүснэгт~\ref{tab:mcnemar}-т үзүүлэв.

\begin{table}[H]
\centering\small
\caption{Test-ийн үр дүн}
\label{tab:test}
\begin{adjustbox}{max width=\textwidth}
\begin{tabular}{llrrrrl}
\toprule
ID & Тохиргоо & Нийт /200 & Хариулах /151 & Боломжгүй /24 & Хоёрдмол /25 & GPU \\
\midrule
D2 & LoRA + RAG & 42 (21.0\%) & 37 (24.5\%) & 0 & 5 & A100 \\
H4 & + засвар, орчуулга & 68 (34.0\%) & 48 (31.8\%) & 9 & 11 & A100 \\
E2 & + 8 хувилбар, санал & 77 (38.5\%) & 62 (41.1\%) & 5 & 10 & T4 \\
E3c & + тодруулга & \textbf{93 (46.5\%)} & 59 (39.1\%) & 17 & 17 & GPU-гүй \\
E4 & + шүүлт + dev-ийн дүрэм & 90 (45.0\%) & --- & --- & --- & GPU-гүй \\
I3t & H4 + 579 жишээ + засварласан орчуулга & 77 (38.5\%) & 62 (41.1\%) & --- & --- & T4 \\
I3tr & I3t + E4-ийн дүрэм & 90 (45.0\%) & --- & --- & --- & GPU-гүй \\
\bottomrule
\end{tabular}
\end{adjustbox}
\end{table}

\begin{table}[H]
\centering\small
\caption{McNemar тест (exact)}
\label{tab:mcnemar}
\begin{tabularx}{\textwidth}{X p{3.0cm} p{2.4cm} r}
\toprule
Харьцуулалт & Хэмжүүр & Ахисан / буурсан & $p$ \\
\midrule
dev D2 → H4 & нийт & +77 / −14 & $10^{-11}$ \\
dev H4 → E1 & нийт & +8 / −7 & 1.0 \\
dev H4 → E2 & нийт & +27 / −14 & 0.06 \\
dev H4 → I3t & хариулах боломжтой & +33 / −10 & 0.0006 \\
test D2 → H4 & нийт & +42 / −16 & $8.6 \times 10^{-4}$ \\
test D2 → E2 & нийт & +48 / −13 & $7.7 \times 10^{-6}$ \\
test H4 → E2 & хариулах боломжтой & +17 / −3 & 0.0026 \\
test E2 → E4 & нийт & +14 / −1 & 0.001 \\
\bottomrule
\end{tabularx}
\end{table}

\textbf{Учир шалтгаан.} Суурь загварт ижил асуултын англи хувилбар кирилл хувилбараас 10 дахин олон зөв хариулт өгсөн
(43 ба 4) тул орчуулга хэлний зөрүүг арилгадаг; хайлтыг монголоор хийх нь (H4) монгол тайлбар, утгатай мэдлэгийн санг
бүрэн ашиглах боломж олгосон. Нэмэлт сургалт бага нөлөөтэй байсны шалтгаан нь D2-ийн сургалтаас өөрийн DB-ийн 344 жишээ
бүгд prompt-ийн уртаас болж хасагдсан явдал --- аргын бус, өгөгдлийн хязгаарлалт. 8 хувилбарын дээд хязгаар (pass@8)
хариулах боломжтой асуултад 127/224 (56.7\%) байхад санал хураалт 100-г сонгосон. H4 dev дээр 36.7\%, test дээр 34.0\%
байгаа нь dev-д хэт тааруулсан шинж бага; E4-ийн дүрэм dev дээр +12.0, test дээр +6.5 пункт нэмсэн.

\textbf{Алдааны шинжилгээ.}\label{sec:errors} E2-ийн dev дээрх хариулах боломжтой 124 буруу асуултын шалтгааныг ангилав
(Хүснэгт~\ref{tab:errors}). Орчуулгын асуудалтай асуултын зөв нь 23.1\%, бусдынх 51.2\%; орчуулгагүй H1-д эдгээр нь
21.2\% ба 20.9\% байсан --- асуулт өөрөө хэцүү биш, орчуулга буруутай. D2-ийн гүйцэтгэлийн 47 алдааны 31 нь boolean
баганыг \tsql{= 0/1}-ээр харьцуулсан алдаа байсан бөгөөд диалект хөрвүүлэлтэд засагдсан.

\begin{table}[H]
\centering\small
\caption{E2-ийн алдааны ангилал (dev, хариулах боломжтой 124 буруу)}
\label{tab:errors}
\begin{tabularx}{\textwidth}{X r}
\toprule
Шалтгаан & Тоо \\
\midrule
Орчуулга --- маскын алдаа 24 («байсан» → салбарын нэр «Баасан»), нэр томьёо 23 («төгрөгийн» → Tenge) & \textbf{47} \\
Нэр томьёо, KPI-ийн тодорхойлолтыг хэрэглээгүй & 22 \\
Хүснэгт / холбоос / багана & 19 \\
Нэгтгэл / томьёо / логик & 18 \\
Утга / литерал & 9 \\
Огноо, гаралтын багана, валют г.м. & 9 \\
\bottomrule
\end{tabularx}
\end{table}

Орчуулгын алдааны бодит жишээ (dev K035): «2026 оны 2-р улиралд \emph{барагдаж хаагдсан} кейсүүдээс нийт хэдэн төгрөг
цуглуулсан бэ? \emph{Долларын} зээлийн кейсийг …» асуултыг орчуулагч «… collected from \emph{closed cases} … convert the
\emph{"Долларын"} loan cases to \emph{Tenge} …» гэж орчуулсан. «Барагдаж» гэсэн утга алдагдсан тул загвар
\tsql{status = 'CLOSED'} гэж шүүсэн (зөв нь \tsql{'RESOLVED'}); «Долларын» орчуулагдалгүй үлдэж, «төгрөг» тенге болсон. Ийм
алдааг багасгахын тулд орчуулгын зааварт мэдлэгийн сангийн нэр томьёог нэмж (T2), цаашлаад санхүүгийн хэллэгээр
орчуулагчийг нэмж сургаж байна (§\ref{sec:mt-train}).

\textbf{Домэйн ба хүндрэлээр.} Хүснэгт~\ref{tab:by-domain}-д хариулах боломжтой асуултын зөвийг домэйн ба хүндрэлээр
задлав. Бүх домэйнд ахиц гарсан ч домэйн хоорондын (cross), эргэн төлөлт (payments), ерөнхий дэвтрийн хандлага (gl\_trends)
асуулт сул хэвээр бөгөөд эдгээр нь олон хүснэгтийн холбоос, хугацааны цуваа, KPI-ийн томьёо шаарддаг.

\begin{table}[H]
\centering\small
\caption{Хариулах боломжтой асуултын зөв, домэйн ба хүндрэлээр (dev)}
\label{tab:by-domain}
\begin{tabularx}{\textwidth}{X r r}
\toprule
Домэйн / хүндрэл & D2 & E2 \\
\midrule
Зээлийн хүсэлт (applications) & 6/28 & 15/28 \\
Барагдуулалт (collection) & 4/28 & 8/28 \\
Домэйн хоорондын (cross) & 0/28 & 5/28 \\
Харилцагч (customers) & 10/28 & 17/28 \\
Хадгаламж (deposits) & 5/28 & 17/28 \\
Ерөнхий дэвтэр, хандлага (gl\_trends) & 4/27 & 9/27 \\
Зээл (loans) & 7/29 & 20/29 \\
Эргэн төлөлт (payments) & 1/28 & 9/28 \\
\midrule
Хялбар / дунд / хүнд & 16/57 · 21/121 · 0/46 & 37/57 · 50/121 · 13/46 \\
\bottomrule
\end{tabularx}
\end{table}

\section{Орчуулгын загварын үр дүн}
\label{sec:mt-results}

Гурван орчуулагчийн dev (300 асуулт) дээрх үзүүлэлтийг Хүснэгт~\ref{tab:mt-results}-д үзүүлэв. MADLAD-ийн латин
асуултыг эхлээд кирилл рүү хөрвүүлж (сайжруулсан галиглал) орчуулсан. Нэмж сургасан MADLAD (M1) орчуулгын бүх хэмжүүрээр
хамгийн сайн гарсан: chrF++ кирилл 74.8, латин 73.6, санхүүгийн нэр томьёоны таарц 69.0\% нь dev-ийн жишиг англи
орчуулгын таарцаас (59.6\%) ч өндөр. Dev-ийн 300 асуулт сургалтын 8 500 эх өгүүлбэртэй яг давхцаагүйг шалгасан. Харин
DB-ийн утгыг (код, шошго) өөрчлөхгүй хадгалсан хувь M1-д 13.6\% (n = 22) буюу Qwen-ийнхоос (63.6\%) муу байна.

Орчуулгын chrF++ өндөр байх нь Text2SQL-ийн үр дүн сайн гэсэн үг биш: zero-shot MADLAD (M0)-ийн chrF++ Qwen-ийнхтой ойролцоо
(53.5 ба 58.0) байсан ч түүгээр орчуулсан асуултын dev EX (I3 тохиргоо, хариулах боломжтой 224) 36.2\% буюу Qwen-ийн
орчуулгатай I3t-ээс (46.9\%) 10.7 пунктээр бага гарсан. Харин нэмж сургасан M1-ээр орчуулсан ижил тохиргоо (M1a) хариулах
боломжтой 127/224 буюу \textbf{56.7\%} (95\% итгэх завсар 50.5--63.4\%), нийт 160/300 (53.3\%) зөв хариулж, I3t-ээс
22 асуултаар (+9.8 пункт) илүү гарав. Энэ нь нэг үүсгэлттэй, dev-ийн дүрэмгүй тохиргоогоор E4-тэй (159/300) тэнцэх үр дүн
юм. Гэвч үр дүн зөвхөн dev дээр, нэг seed-ээр хэмжигдсэн тул test дээр нэг удаа баталгаажуулах шаардлагатай.

\begin{table}[H]
\centering\small
\caption{Орчуулгын үзүүлэлт (dev, 300 асуулт)}
\label{tab:mt-results}
\begin{tabularx}{\textwidth}{X r r}
\toprule
Үзүүлэлт & Кирилл & Латин \\
\midrule
chrF++ (Qwen3.5-9B, засварласан кэш) & 58.0 & 33.9 \\
chrF++ (MADLAD-400-3B zero-shot, M0) & 53.5 & 51.1 \\
chrF++ (MADLAD-400-3B QLoRA, M1) & \textbf{74.8} & \textbf{73.6} \\
Нэр томьёо зөв (Qwen / M1) & 46.3\% / \textbf{69.0\%} & 26.7\% / \textbf{67.5\%} \\
Text2SQL EX, хариулах боломжтой (Qwen I3t / M0a / M1a) & \multicolumn{2}{r}{46.9\% / 36.2\% / \textbf{56.7\%}} \\
Латин → кирилл хөрвүүлэлтийн үгийн нарийвчлал & --- & 89.1\% → 94.8\% \\
Маскласан асуулт (T2 засвар, dev+test 500) & \multicolumn{2}{r}{122 → 37} \\
Орчуулгын нөлөө Text2SQL-д (D2 → D2 tr, dev нийт зөв) & \multicolumn{2}{r}{+11.3 пункт} \\
\bottomrule
\end{tabularx}
\end{table}

\section{Системийн хугацааны хоцрогдлын үнэлгээ}
\label{sec:latency}

\begin{table}[H]
\centering\small
\caption{Тохиргооны асуулт тутмын дундаж хугацаа (Kaggle T4×2)}
\label{tab:latency}
\begin{tabularx}{\textwidth}{X r p{4.2cm}}
\toprule
Тохиргоо & с / асуулт & Тайлбар \\
\midrule
H4, I3t (нэг үүсгэлт, орчуулгатай) & ~20--30 & Нийт dev 300 ≈ 2--2.5 цаг \\
B бүлэг (урт RAG prompt, анхны хувилбар) & ~53 & Prompt 13--20 мянган токен \\
E2 (8 хувилбар + засвар) & ~175 & dev 300 ≈ 14.6 цаг \\
E3c, E4 (дахин сонголт) & ≈ 0 & GPU-гүй, E2-ийн хувилбараас \\
\bottomrule
\end{tabularx}
\end{table}

Нэг үүсгэлттэй I3t test дээр 38.5\% буюу E2-тэй тэнцүү гарсан нь 6--8 дахин бага тооцооллоор ижил үр дүнд хүрэх
боломжийг харуулж байна; бодит хэрэглээнд энэ тохиргоо илүү тохиромжтой.

\section{Зорилго ба таамаглалын биелэлт}
\label{sec:hyp-check}

\begin{table}[H]
\centering\small
\caption{Зорилго ба таамаглалын биелэлт}
\label{tab:hyp-check}
\begin{tabularx}{\textwidth}{p{1.0cm}X p{2.4cm}}
\toprule
№ & Нотолгоо & Дүгнэлт \\
\midrule
H1 & A7 → B5 +12.0 пункт (dev); H4 → I3t (жишээний сан 160 → 579) +33/−10, $p = 0.0006$ & Батлагдсан \\
H2 & D2 → D2 tr +11.3 пункт; test D2 → H4 $p = 8.6\times10^{-4}$; сургасан орчуулагч I3t → M1a +9.8 пункт (dev) & Батлагдсан \\
H3 & B5 → D2 +1.4 пункт; өөрийн DB-ийн жишээ сургалтад ороогүй & Батлагдаагүй (D3-аар шалгана) \\
H4 & test H4 → E2 хариулах боломжтой +17/−3, $p = 0.0026$; нийт $p = 0.12$ & Хэсэгчлэн батлагдсан \\
\bottomrule
\end{tabularx}
\end{table}

\textbf{Үр дүнгийн хязгаарлалт.}\label{sec:limitations}\label{sec:deviations}
(1) E4-ийн татгалзал, тодруулгын түлхүүр үг ба E3c-ийн 0.5 босгыг dev-ээс гаргасан тул test-ийн тоог гол гэж үзнэ.
(2) Асуулт ба gold SQL-ийг LLM agent-ууд бичиж, өөр agent-ууд хянасан; хүний шалгалт хийгдээгүй. (3) Test-ийн D2, H4 нь
A100, E2, E3c, I3 нь T4 дээр хийгдсэн тул H4 → E2 харьцуулалтад төхөөрөмжийн нөлөө холилдсон. (4) Test 200 асуулт
(±7 пункт орчим). (5) E3c dev-д 300-аас 110 асуултад тодруулга асуусан нь бодит хэрэглээнд хэт олон. (6) Төлөвлөгөөний
B6--B10, C4--C7 (GRPO), F, G (agentic) бүлэг GPU-ийн төсвийн улмаас ажиллаагүй; код нь бэлэн.

\section{Ёс зүйн асуудал}
\label{sec:ethics}

\textbf{Өгөгдлийн зөвшөөрөл ба нууцлал.} Бодит харилцагчийн мэдээлэл огт ашиглаагүй; өгөгдлийн сан бүхэлдээ синтетик.
Бодит байгууллагад ашиглагдах тохиолдолд хувийн мэдээллийн баганыг хассан давхарга, дүрмээр татгалзах, унших эрх ба RLS,
шифрлэгдсэн аудит, байгууллагын дотор ажиллах нээлттэй эхийн загвар гэсэн давхаргат хамгаалалт хийв.

\textbf{Хазайлт ба хортой хэрэглээ.} Суурь загвар кирилл асуултад 1.3\%, латин асуултад 0\% зөв хариулсан нь латинаар
бичдэг хэрэглэгчдэд чанар муутай үйлчилгээ үзүүлэх эрсдэлтэй тул бичиг тус бүрээр тусад нь хэмжих шаардлагатай. Буруу
тоонд үндэслэсэн шийдвэр хохирол учруулж болох тул систем итгэл багатай үед тодруулга асууж, ашигласан хүснэгт ба SQL-ийг
харуулдаг; SQL ба prompt injection-ийг шалгаж хаасан. Систем нь шийдвэр гаргагчийг орлох бус, туслах хэрэгсэл юм.

\textbf{Хиймэл оюуны хэрэгслийн ашиглалт.}\label{sec:ai-tools} Системийн кодыг LLM кодын туслахын (Claude Code) тусламжтайгаар
бичиж, бие даасан «tester» agent-аар логикийн алдааг шалгуулсан. Үнэлгээний асуулт, gold SQL, жишээ, орчуулгын хосыг LLM
agent-ууд ноороглож, өөр agent-ууд сохроор хянасан; gold SQL бүрийг DB дээр гүйцэтгэж шалгасан. Иймд үнэлгээний санг
«хүний бичсэн жишиг» гэж нэрлээгүй.

\section{Бүлгийн дүгнэлт}
\label{sec:ch7-summary}

Test дээр систем нэмэлт сургалттай суурийн 21.0\%-иас 45--46.5\% хүрч, ахиц нь статистикийн хувьд ач холбогдолтой. RAG ба
орчуулга хамгийн их нэмэр өгсөн, санал хураалт хариулах боломжтой асуулт дээр ач холбогдолтой нэмэр өгсөн. Алдааны 38\%
нь орчуулгаас үүдсэн тул орчуулагчийг сайжруулах нь хамгийн том нөөц бөгөөд латин асуултын орчуулга хамгийн сул хэвээр.
